% DRAFT — 并入 main.tex 的新小节(T 行为半区数字留槽)
% 位置:插在 5.5 (Determinism re-run) 之后、5.6 Additional observations 之前;
% 后续小节编号顺延。数字全部程序化提取自 capability_gap_v2 归档。

\subsection{Comparative behavioral evidence (campaign v2)}
\label{sec:res-campaign}

The protocol above tests activation against closed-form and embedding
retrieval on \emph{mechanism} metrics. To test the same questions on
\emph{behavior}, we ran a second campaign (360 runs; design, two
mid-course confound repairs, and full data in Supplementary
Material~S15--S17) in which every condition shares the same scripted
experience stream, the same menu, the same LLM decision head
(temperature 0), and the same decision budget---only the memory form
differs.

\textbf{Goal resumption under neutral instructions (C).} After an
open-book practice phase (obtain \texttt{oak\_planks}) and an unrelated
distractor phase, the test prompt gives no goal at all
(``you have free choice'') and the recipe facts are stripped. FAS
resumes and completes the abandoned goal in \textbf{9/10} paired seeds;
the full-history LLM baseline and the no-memory baseline resume in
\textbf{0/10} (McNemar $p{=}0.0039$; Holm $p{=}0.0117$). This advantage
is regime-dependent: with a longer experience stream that also contains
a failure episode, the full-history baseline matches FAS at 10/10
(below)---the salience advantage holds in the short-history regime.
Mechanism attribution: a zero-experience control reproduces the
same working set (the recipe-closure prior edge
\texttt{oak\_log}$\rightarrow$\texttt{oak\_planks} re-lights the chain
when the raw material is observed), so the effect is
\emph{observation-triggered activation over prior relational
structure}---relational working-set salience---rather than an
intention-lifecycle mechanism.

\textbf{Goal revision is where the same mechanism costs (E).} With an
infeasible target (tool-gated ore), the full-history baseline stops
repeating the failed gather after \textbf{1.0} failures on average and
pivots immediately; FAS repeats it \textbf{10.5} times and pivots in
7/10 runs (late). A joint measurement (below) gives the direction and
magnitude of this cost precisely. The two results are the same coin: state that keeps
a goal salient under distraction also keeps it salient when the
environment invalidates it.

\textbf{Consumer repairs (E2/F).} We then tested whether the failure
and goal states are simply not consumed by the decision context.
Wiring the production event-frame's negative-polarity encoding into
the campaign assembly (failure$\rightarrow$entity edge $w{=}{-}0.8$,
parameters identical to the production config) reduced repeated
failures from 10.5 to \textbf{7.2} (one-sided Wilcoxon
$p{=}0.011$)---statistically significant, still short of the baseline,
so residual inertia is design-level (constant goal text, weak ledger
prior). Adding a uniform self-graph goal consumer
(\texttt{current\_goal} $\rightarrow$ working set; production default
off) turned the previously bit-identical ON/OFF behavior into a
separation: with the registry cleared, the consumer-on loop re-derives
the abandoned goal in 5/5 seeds and its gathering actions rise from
1.0 to \textbf{6.0} per run. The self-model's ``behavioral
irrelevance'' was a missing-consumer artifact, not a property of the
representation.

\textbf{Behavioral convergence (I).} On a two-chain convergence task
(\texttt{oak\_fence} requires \texttt{oak\_planks} and \texttt{stick}
learned in separate experiences), FAS conditions complete 0--2/10
while the raw-history LLM baseline completes \textbf{10/10} (Holm
$p{=}0.041$). Context budget ($k{=}16,32$) and a structure-enriched
serialization (explicit action$\rightarrow$outcome pairs) both remain
at 0/10 ($n{=}30$), locating the failure in the \emph{serialization
format}: the graph state drops the temporal narrative that the raw
stream preserves (Supplementary Material~S17).

\textbf{Credit horizon (B) and long-run stability (J).} A deterministic
delay probe on the production attribution pipeline shows
action$\rightarrow$consequence credit forms in 10/10 seeds up to 80\,s
and in 0/10 at 160\,s---the 90\,s attribution window is the
architecture's credit horizon (an intentional semantic boundary,
Section~\ref{sec:limits}). A 100-cycle stability run (zero-LLM) shows
bounded graph growth ($+3.7\%$ nodes), clamped activation, and no
contamination collapse, but also zero promotion (repetition-starved
episodes never reach the support threshold)---long-run consolidation
remains untested.

\textbf{Joint measurement (T).} On a single task sequence carrying
both mechanisms (a failed infeasible goal followed by a neutral free
phase), the event-frame encoding replicates its effect
(failed-gather repetitions 12.0$\rightarrow$7.8, one-sided Wilcoxon
$p{=}0.002$), the goal consumer adds no resumption increment where the
prior-edge path is already at ceiling (all conditions 10/10) and adds
no inertia penalty (6.8 vs 7.8, n.s.)---the two consumers coexist, and
the combined working set crowds the failed action out of the top-8
entirely. The behavioral trade-off curve between suppression and
salience is therefore measurable, and its first two points favor
suppression. (One design flaw---inventory carry-over saturating the
first resumption metric---was caught and corrected with a fresh-world
rerun; both datasets are archived.)

Together these results revise the paper's claim structure: persistent
relational state demonstrably changes behavior in one direction
(resumption), demonstrably costs in another (revision), is currently
bottlenecked at the serialization interface (convergence), and is
credit-limited by an explicit design window (delayed consequence).
All negative and reverse-direction results are reported at full
strength.
